\section{Sistematika Penulisan}

Tesis ini disusun dalam lima bab utama.

Bab I memuat pendahuluan yang mencakup latar belakang masalah, rumusan masalah, tujuan penelitian, kontribusi penelitian, hipotesis dan premis, batasan penelitian, metodologi penelitian, serta sistematika penulisan.

Bab II memuat tinjauan pustaka dan landasan teori mengenai IoT, keamanan dan kontrol akses pada IoT, ABAC, arsitektur \textit{edge/fog computing, attribute freshness}, mekanisme \textit{cache} atribut, TTL, algoritma metaheuristik, serta posisi penelitian ini di antara studi-studi lainnya.

Bab II memuat metodologi penelitian yang menjelaskan model konseptual, arsitektur ABAC berbasis \textit{edge, threat model}, desain mekanisme \textit{cache} atribut, formulasi optimasi \textit{cacheable attribute}, TTL, dan ukuran \textit{cache}, desain algoritma \textit{GWO-optimized ABAC}, skenario \textit{workload, baseline} eksperimen, metrik evaluasi, serta prosedur pengujian.

Bab IV memuat implementasi, pengujian, analisis, dan pembahasan, termasuk detail implementasi simulasi berbasis iFogSim2 dan Java, konfigurasi lingkungan eksperimen, hasil pengujian pada berbagai skenario dinamika, serta perbandingan kinerja antar konfigurasi uji. Bab ini juga membahas \textit{trade-off} antara latensi, \textit{PIP calls, cache efficiency, attribute freshness, false allow rate, false deny rate, stale cache hit rate,} dan \textit{revocation false allow rate}.

Bab V memuat kesimpulan dan saran untuk penelitian lanjutan. Bab ini merangkum temuan utama penelitian, menilai pencapaian tujuan penelitian, menjelaskan kontribusi \textit{GWO-optimized ABAC} terhadap efisiensi evaluasi akses dan kualitas keputusan otorisasi, serta memberikan arah pengembangan lebih lanjut, seperti pengujian pada \textit{testbed} IoT fisik, integrasi dengan mekanisme autentikasi nyata, dan implementasi \textit{event-based simulation} yang lebih lengkap pada iFogSim2.